\vssub
\subsection{~数据同化接口} \label{sec:das}
\vssub

如上所述，波浪模式子程序还配有一个数据同化接口例程（{\file w3wdasmd.ftn} 中的 {\F w3wdas}）。该例程集成在独立壳中（见 \para\ref{sec:ww3shel}），用于为组合的波浪模式/数据同化方案提供时间步管理。它尚未集成到多网格模式驱动器中，尽管多网格模式管理算法已经考虑了它。

这里假定采用一种相当简单的方法：数据同化在选定时刻执行，而波浪模式持续向前积分。在程序壳的设置中，数据同化在模式到达目标时间之后、尚未产生输出之前执行。完成数据同化后，只为按请求生成输出而再次调用波浪模式例程。因此，给定时刻的波浪模式输出将包含该目标时刻数据同化的影响。

通用程序壳还处理若干类型的待同化数据，并将其传递给数据同化接口例程。所有数据都需要使用波浪模式输入预处理器进行预处理（见 \para\ref{sec:ww3prep}），并由通用壳通过文件名识别。目前最多可使用三个不同的数据文件。暂定而言，它们可分别为平均波浪参数、一维谱数据和二维谱数据。不过，这并未硬编码到模式中，实际上需要由用户定义。

目前尚无可用的数据同化软件包。因此，需要用户提供数据同化方案，并可使用接口例程（{\file w3wdasmd.ftn} 中的 {\F w3wdas}）将其纳入波浪模式；该例程的文档应足以支持必要的编程。关于如何将用户提供的软件添加到 \ws{} 编译系统，详见下一章。虽然 NCEP 目前正在研究波浪数据同化技术，但没有计划将波浪数据同化软件作为 \ws{} 软件包的一部分发布。
